\chapter{Lo que no se sabe y lo que salió mal}
\label{cap:s-honestidad}
Un proyecto de datos convence más cuando dice con claridad dónde termina lo que sabe. Este capítulo junta los huecos, los
supuestos y los errores que se corrigieron. Es material para el coloquio: si un profesor encuentra un límite, lo mejor es
que ya esté escrito aquí.

\section{Lo que no se pudo medir (huecos)}
\begin{center}
\small
\begin{tabularx}{\linewidth}{>{\raggedright\arraybackslash}p{0.38\linewidth}Y}
\encabezado \h{Qué falta} & \h{Qué se hizo en su lugar} \\
Ocupación hotelera del sur desde 2025 & Se mide con lo que sí llega a 2026 (Tren Maya, Belice, INAH) y se muestra ``sin
dato oficial'' \\
\filagris Visitantes de Maya Ka'an y de la Laguna Milagros & No se pronostican; en el Radar, ``sin dato oficial'' \\
De qué estado viene el turista mexicano, su edad y su ingreso & Se marca como hueco en la viajera ideal \\
\filagris Conversión de Facebook para turismo & Se prueba con tres valores (3\,\%, 5.75\,\% y 6.38\,\%) \\
Reseñas de los cinco lugares & Se usan las de los lugares llenos para saber qué molesta; en la página, un botón a Google \\
\filagris Fotos de platillos tomadas en el sur & Se muestran fotos reales de restaurantes; las de platillos las aportará el
equipo \\
Estrellas oficiales de los hoteles del sur & Solo existen para 70 centros del país, ninguno del sur \\
\filagris Tormentas de 2026 & ``Sin dato'', nunca ``sin tormenta'' \\
Prueba de la página con personas & Se declara que no se hizo \\
\bottomrule
\end{tabularx}
\normalsize
\end{center}

\section{Los supuestos (y cómo se probaron)}
Un supuesto es un número que no se pudo medir y se tuvo que fijar. Cada uno está marcado y se probó con varios valores:
\begin{itemize}
\item \textbf{Presupuesto de \$250,000 al año}: no hay un presupuesto real. Se probó con \$500,000 y \$1,000,000: los
  visitantes crecen en proporción y el reparto casi no cambia.
\item \textbf{Costos de anuncios}: son promedios de anunciantes de Estados Unidos.
\item \textbf{Cuánto quita una tormenta}: 0\,\%, 25\,\% o 50\,\%; el año casi no cambia.
\item \textbf{Cada viaje planeado es un visitante más}: es una cota alta; si llegan menos, la capacidad queda aún más
  protegida.
\end{itemize}

\section{Los límites de los modelos}
\begin{itemize}
\item El Radar apenas le gana a ``igual que este mes'' (130 contra 126 de 156) y en los cinco lugares casi no hubo cambios
  con qué probarlo.
\item El rango del 90\,\% de la Ruta se cumple 80 de cada 100 veces, no 90, por la caída inesperada de 2023.
\item El clima casi no mejora el pronóstico.
\item Markov y la ocupación semanal solo cubren el norte.
\item SECTUR mide menos ocupación que el gobierno del estado (hasta 19 puntos en Isla Mujeres). Como el norte se ve
  \emph{menos} lleno de lo que diría el estado, el error juega en contra de la conclusión ``el sur tiene espacio'', no a
  favor: la conclusión se sostiene igual.
\end{itemize}

\section{Los errores que se encontraron y se corrigieron}
Estos errores son parte del valor del proyecto: muestran que las cifras se revisaron y no se aceptaron a ciegas.

\begin{enumerate}
\item \textbf{El conteo del DENUE} sumaba los diccionarios de datos como negocios (1,914 de más). Lo detectó Spark al
  contar diferente; un tercer lector lo confirmó.
\item \textbf{El Censo} contaba un catálogo como pueblos: 2,257 en lugar de 2,243.
\item \textbf{La ``ocupación'' de 2,560,068}: el panel del Radar sumaba tres renglones distintos (cuartos disponibles,
  ocupados y porcentaje). Ahora se calcula ocupados entre disponibles, y una prueba lo vigila.
\item \textbf{Visitantes contados dos veces}: la serie de ``zonas arqueológicas'' del gobierno del estado es el mismo dato
  del INAH sumado. Se usa solo el INAH.
\item \textbf{Cancún ``tranquilo'' en julio}: el índice falló la prueba de validez y se corrigió con datos (capítulo~5).
\item \textbf{Afluencia que parecía llegar a junio de 2024}: unos ceros lo hacían ver así; en realidad termina en marzo.
\item \textbf{Una línea del archivo de huracanes sin hemisferio}: la primera exploración la saltaba sin avisar. El código
  final se detiene si no entiende una línea, y esa latitud se dejó vacía, sin adivinar.
\item \textbf{Clima raro en el 47\,\% de las semanas}: se cambió el corte (capítulo~9).
\item \textbf{Todo el dinero de la Ruta en junio}: se agregó el desempate proporcional (capítulo~8).
\item \textbf{Junio como mes recomendado}: se agregó la condición de lluvia (capítulo~11).
\item \textbf{Fotos de Mérida y Campeche}: se retiraron y se exigió coordenada comprobada (capítulo~11).
\item \textbf{Un dominio inventado} en las maquetas de los anuncios: se cambió por la dirección real de la página.
\item \textbf{La página tenía 67 fallas de accesibilidad}: se corrigieron hasta 0.
\item \textbf{Una imagen borrada}: la portada del documento ejecutivo desapareció del disco; se restauró y se agregó una
  prueba que revisa que todas las imágenes existan.
\end{enumerate}
